% Minimal retained preamble for lane-p2 cards taken from
% sources/logic-fl05-ramsey-like-theorems-for-the-schre/pinned-original-source.tex
% (arXiv:2412.11598v3, CC BY 4.0).  This source ships no separate preamble file:
% the pinned tex is a concatenation of the authors' main.tex and macros.tex and
% carries its preamble inline, so the blocks below are sliced mechanically from
% that file and nothing has been renamed or redefined.
%
% Copied verbatim from the source:
%   line 2        : the authors' class line
%   lines 3-5,7,9 : graphicx, amsthm, the ams bundle, cleveref, url
%   line 10       : tikz-cd, used by the authors' commutative diagrams
%   lines 20-108  : the authors' macro definitions, taken in full so that every
%                   macro of the retained spans is available
%   lines 113-124 : the authors' theorem declarations
%   lines 131-141 : the authors' machinery for restating a theorem
%   lines 143-148 : the authors' coloured comment macros
%
% Omitted source preamble lines and why:
%   line 1, 110-111 and 150 (concatenation comment, title, date, author): the
%       assembler supplies the title block, author list and date;
%   line 6 (hyperref) and line 8 (inputenc with the utf8x option): the forced
%       xelatex engine supplies the hyperref wiring itself and utf8x is
%       incompatible with it;
%   line 115 (numberwithin theorem by section): the assembler numbers the
%       retained material itself;
%   lines 120 and 121-123 are declarations that the authors left commented out or
%       that duplicate earlier ones and are not reproduced.
\documentclass{article}
\usepackage{graphicx}
\usepackage{amsthm}
\usepackage{amsfonts, amsmath, amssymb, amsthm}
\usepackage{cleveref}
\usepackage{url}
\usepackage{tikz-cd}
\newcommand{\Psf}{\mathsf{P}}
\newcommand{\Qsf}{\mathsf{Q}}
\newcommand{\Rsf}{\mathsf{R}}
\def\F{\mathcal{F}}
\def\U{\mathcal{U}}
\def\C{\mathcal{C}}
\def\I{\mathcal{I}}
\def\D{\mathcal{D}}
\def\L{\mathcal{L}}
\def\V{\mathcal{V}}
\def\A{\mathcal{A}}
\def\B{\mathcal{B}}
\def\P{\mathcal{P}}
\def\M{\mathcal{M}}
\def\Nc{\mathcal{N}}
\def\Q{\mathcal{Q}}
\newcommand{\cs}{2^{\NN}}
\newcommand{\bs}{\NN^{\NN}}
\newcommand{\bstr}{2^{<\NN}}
\newcommand{\uh}{\upharpoonright}
\newcommand{\converge}{\!\!\downarrow}
\newcommand{\diverge}{\!\!\uparrow}
\newcommand{\qvdash}{\operatorname{{?}{\vdash}}}
\newcommand{\nqvdash}{\operatorname{{?}{\nvdash}}}
\newcommand{\ISig}{\mathsf{I}\Sigma^0}
\newcommand{\Robinson}{\mathsf{Q}}
\newcommand{\IPi}{\mathsf{I}\Pi^0}
\newcommand{\IDelta}{\mathsf{I}\Delta^0}
\newcommand{\Iopen}{\mathsf{I_{open}}}
\newcommand{\BSig}{\mathsf{B}\Sigma^0}
\newcommand{\BPi}{\mathsf{B}\Pi^0}
\newcommand{\BDelta}{\mathsf{B}\Delta^0}
\newcommand{\BCSig}{\mathsf{BC}\Sigma^0}
\newcommand{\BCPi}{\mathsf{BC}\Pi^0}
\newcommand{\BCDelta}{\mathsf{BC}\Delta^0}
\newcommand{\LSig}{\mathsf{L}\Sigma^0}
\newcommand{\LPi}{\mathsf{L}\Pi^0}
\newcommand{\CSig}{\mathsf{C}\Sigma}
\newcommand{\RCA}[0]{\mathsf{RCA}}
\newcommand{\WKL}[0]{\mathsf{WKL}}
\newcommand{\WWKL}[0]{\mathsf{WWKL}}
\newcommand{\ACA}[0]{\mathsf{ACA}}
\newcommand{\ATR}[0]{\mathsf{ATR}}
\newcommand{\PIOOCA}[0]{\Pi^1_1\mathsf{-CA}}
\newcommand{\CAC}[0]{\mathsf{CAC}}
\newcommand{\ADS}[0]{\mathsf{ADS}}
\newcommand{\RRT}[0]{\mathsf{RRT}}
\newcommand{\FS}[0]{\mathsf{FS}}
\newcommand{\PFS}[0]{\mathsf{PFS}}
\newcommand{\TS}[0]{\mathsf{TS}}
\newcommand{\GRRT}[0]{\mathsf{GRRT}}
\newcommand{\LRRT}[0]{\mathsf{LRRT}}
\newcommand{\EM}[0]{\mathsf{EM}}
\newcommand{\SEM}[0]{\mathsf{SEM}}
\newcommand{\HalfEM}[0]{\mathsf{Half}\mbox{-}\mathsf{EM}}
\newcommand{\RT}[0]{\mathsf{RT}}
\newcommand{\SRT}[0]{\mathsf{SRT}}
\newcommand{\NN}[0]{\mathbb{N}}
\newcommand{\nat}[0]{\mathbb{N}}
\newcommand{\card}{\operatorname{card}}
\newcommand{\finsub}{\subseteq_{\mathtt{fin}}}
\newcommand{\W}{\mathrm{W}}
\newcommand{\sW}{\mathrm{sW}}
\newcommand{\lex}{\mathrm{lex}}
\newcommand{\base}{\operatorname{base}}
\newcommand{\range}{\operatorname{ran}}
\def\qt#1{``#1''}%

\newtheorem*{statement}{Statement}
\newtheorem{theorem}{Theorem}
\numberwithin{theorem}{section}
\newtheorem{maintheorem}[theorem]{Main Theorem}
\newtheorem{lemma}[theorem]{Lemma}
\newtheorem{question}[theorem]{Question}
\newtheorem{proposition}[theorem]{Proposition}
%\newtheorem{statement}[theorem]{Statement}
\newtheorem{remark}[theorem]{Remark}
\newtheorem{definition}[theorem]{Definition}
\newtheorem{corollary}[theorem]{Corollary}
\newtheorem{example}[theorem]{Example}

% Allow to repeat theorems
\makeatletter
\newtheorem*{rep@theorem}{\rep@title}
\newcommand{\newreptheorem}[2]{%
\newenvironment{rep#1}[1]{%
 \def\rep@title{#2 \ref{##1}}%
 \begin{rep@theorem}}%
 {\end{rep@theorem}}}
\makeatother

\newreptheorem{theorem}{Theorem}

% Comments
\usepackage{xcolor}
\newcommand{\ludovic}[1]{\textcolor{blue}{#1}}
\newcommand{\quentin}[1]{\textcolor{green}{#1}}
\newcommand{\lorenzo}[1]{\textcolor{orange}{#1}}
\newcommand{\oriola}[1]{\textcolor{red}{#1}}
\crefname{originallemma}{original lemma}{original lemmas}
\Crefname{originallemma}{Original lemma}{Original lemmas}
\makeatletter
\newlabel{lem:fs-stabilizing}{{6.18}{source}{}{sourcepointer-1}{}}
\newlabel{lem:fs-stabilizing@cref}{{[originallemma][0][]6.18}{[1][0][]source}{}{}{}}
\makeatother
\crefname{originallemma}{original lemma}{original lemmas}
\Crefname{originallemma}{Original lemma}{Original lemmas}
\makeatletter
\newlabel{lem:fs-preserving-a}{{6.16}{source}{}{sourcepointer-2}{}}
\newlabel{lem:fs-preserving-a@cref}{{[originallemma][0][]6.16}{[1][0][]source}{}{}{}}
\makeatother
\crefname{originallemma}{original lemma}{original lemmas}
\Crefname{originallemma}{Original lemma}{Original lemmas}
\makeatletter
\newlabel{lem:fs-preserving-b}{{6.20}{source}{}{sourcepointer-3}{}}
\newlabel{lem:fs-preserving-b@cref}{{[originallemma][0][]6.20}{[1][0][]source}{}{}{}}
\makeatother
\title{Every condition extends to one with a strictly longer stem}
\author{Lorenzo Carlucci, Oriola Gjetaj, Quentin Le Houérou, Ludovic Levy Patey}
\date{arxiv:2412.11598v3; CC BY 4.0}
\begin{document}
\maketitle
\noindent\textit{Editorial note.} The counted unit is the article's own Lemma lem:fs-extendibility at source lines 1179-1181 together with its proof at source lines 1182-1185, a straight adjacent pairing in the pinned file. The pinned file is a concatenation of the authors' main and macro files; it is the article's fixed version arXiv:2412.11598v3 under CC BY 4.0.\par\medskip
\noindent\textit{Editorial note.} This source ships no separate preamble file, so the retained preamble is sliced mechanically from the pinned file by the lane and is shared, byte for byte, with the other cards of this source; it carries the authors' full macro block and the cleveref package that the retained proof needs for its capitalised cross-references.\par\medskip
\noindent\textit{Editorial note.} Two context spans carry the material the retained statement and proof rely on: the definition of a condition with its properties (a), (b), (c) and the extension relation at lines 1105-1117, and the definition of the limit colouring at lines 1154-1156. Both are context, not counted units, and neither overlaps the counted interval.\par\medskip
\noindent\textit{Editorial note.} The retained proof names three lemmas of the same section at source lines 1121-1125, 1144-1147 and 1162-1167, each of which is counted by a separate card of this lane. All three are declared as non-load-bearing navigation pointers with the article's own section-wise numbers 6.16, 6.18 and 6.20, counted over the numbered environments of section 6; none of their statements is reproduced inside the span, and no bibliography is needed because the retained spans contain no citation.\par\medskip
\noindent\textit{Editorial note.} The proof leans on the defining property (c) of a condition and on the reservoir inclusion only implicitly, through the word extension; the retained context supplies both, so the card does not depend on material outside its own spans apart from the three pointers.\par\medskip
\noindent\textit{Allowed standard prerequisites (editorial).}\begin{itemize}
\item Conditions, stabilisation, the limit colouring and the three lemmas of the section. A condition is a triple of a colouring constrained by a computable bounding function, a finite stem and an infinite reservoir computing a fixed set D, subject to the defining properties (a) and (b) about how the images of finite sets meet the stem and the reservoir; a reservoir stabilises an interval when the images of its finite sets land in fixed finite sets of short tuples, and the section attaches a limit colouring to such a reservoir. Three lemmas of the same section, stated and proved earlier at source lines 1121-1133, 1144-1150 and 1162-1175, provide respectively a reservoir preserving property (a), a reservoir stabilising the stem, and the preservation of property (b) when a free block is appended to the stem. The author's own definitions and those three lemmas are the allowed inputs.. The counted lemma is the extensibility statement that drives the generic filter construction of the section: starting from any condition it produces one with a strictly longer stem while keeping the defining properties, and the three lemmas are exactly the tools it combines.
\end{itemize}\medskip
\ifdefined\hypertarget\hypertarget{sourcepointer-1}{}\fi
\noindent\textit{Source reference.} Original lemma 6.18: sources/logic-fl05-ramsey-like-theorems-for-the-schre/pinned-original-source.tex, source lines 1144--1147. Navigation pointer invoked by the retained proof to obtain a reservoir that stabilises the interval up to the least element while keeping the cone of D avoided. That lemma is counted by a separate card of this lane, so its proof is not reproduced here. The number is the article's own section-wise number for that lemma, 6.18, counted over the numbered environments of section 6.\par\medskip
\ifdefined\hypertarget\hypertarget{sourcepointer-2}{}\fi
\noindent\textit{Source reference.} Original lemma 6.16: sources/logic-fl05-ramsey-like-theorems-for-the-schre/pinned-original-source.tex, source lines 1121--1125. Navigation pointer invoked by the retained proof to restore property (a) for the enlarged stem by shrinking the reservoir. That lemma is counted by a separate card of this lane, so its proof is not reproduced here. The number is the article's own section-wise number for that lemma, 6.16, counted over the numbered environments of section 6.\par\medskip
\ifdefined\hypertarget\hypertarget{sourcepointer-3}{}\fi
\noindent\textit{Source reference.} Original lemma 6.20: sources/logic-fl05-ramsey-like-theorems-for-the-schre/pinned-original-source.tex, source lines 1162--1167. Navigation pointer invoked by the retained proof to preserve property (b) when the least reservoir element is appended to the stem. That lemma is counted by a separate card of this lane, so its proof is not reproduced here. The number is the article's own section-wise number for that lemma, 6.20, counted over the numbered environments of section 6.\par\medskip
\begin{definition}
A \emph{condition} is a triple $(f, \sigma, X)$ such that
$f \in \P$ is $b_f$-constrained, for some computable function $b_f : \NN \to \NN$,
$\sigma \in [\NN]^{<\omega}$, $X \subseteq \NN$ is an infinite set such that $\max \sigma < \min X$, and
\begin{enumerate}
    \item[(a)] for every~$s \in [\sigma \cup X]^{\leq h(\cdot)}$ with $\min s \in \sigma$, $f(s) \cap X \subseteq s$.
    \item[(b)] for every~$s \in [\sigma \cup X]^{\leq h(\cdot)}$, $f(s) \cap \sigma \subseteq s$.
    \item[(c)] $D \not \leq_T X$.
\end{enumerate}
A condition $d = (g, \tau, Y)$ \emph{extends} $c = (f, \sigma, X)$ (written $d \leq c$) if
$g \leq f$, $\tau \succeq \sigma$, $Y \subseteq X$ and $\tau \setminus \sigma \subseteq X$.
\end{definition}


\begin{definition}
Let $X$ be a reservoir stabilizing $[0, k]$. The \emph{limit coloring} is the function $g_{k,X} : [k]^{\leq h(\cdot)} \to [k]^{<\omega}$ defined by $g_{k,X}(\emptyset) = \emptyset$ and $g_{k,X}(t) = \bigcup_{n \leq h(t) - |t|} I_{t, n}$ otherwise.
\end{definition}

\begin{lemma}\label[lemma]{lem:fs-extendibility}
Let $(f, \sigma, X)$ be a condition. There is an extension $(f, \tau, Y) \leq (f, \sigma, X)$ such that $|\tau| > |\sigma|$.
\end{lemma}

\begin{proof}
Let $x = \min X$. By \Cref{lem:fs-stabilizing}, there is an infinite subset $Y_0 \subseteq X$ stabilizing $[0, x]$ such that $D \not \leq_T Y_0$. Let $g$ be the limit function. By \Cref{lem:fs-preserving-a}, there is an infinite subset $Y \subseteq Y_0$ such that $(f, \sigma \cup \{x\}, Y)$ satisfies property (a).
By \Cref{lem:fs-preserving-b}, $\{x\}$ being vacuously $g$-free, $(f, \sigma \cup \{x\}, Y)$ satisfies property (b). Thus, $(f, \sigma \cup \{x\}, Y)$ is a valid extension.
\end{proof}

\end{document}
